\documentclass[11pt,a4paper]{article}

% ---------------------------------------------------------
% PACKAGES
% ---------------------------------------------------------
\usepackage[utf8]{inputenc}
\usepackage[T1]{fontenc}
\usepackage{lmodern}
\usepackage{amsmath, amssymb, amsfonts, bm}
\usepackage{geometry}
\usepackage{graphicx}
\usepackage{hyperref}
\usepackage{authblk}
\usepackage{cite}
\usepackage{physics}
\usepackage{mathtools}
\usepackage{tensor}
\usepackage{enumitem}

\geometry{margin=2.4cm}

% ---------------------------------------------------------
% TITLE
% ---------------------------------------------------------
\title{\textbf{Companion V (v49.1)}\\[0.2cm]
\textbf{Theoretical C-Coefficients in RDCC}\\[0.2cm]
\large Foundations of Relaxation-Driven Cyclic Cosmology (RDCC)}

\author[1]{Michael Lehmann}
\affil[1]{\small Independent Researcher, Relaxation-Driven Cyclic Cosmology (RDCC)\\
\texttt{mi.lehmann@gmx.de}}

\date{September 2026}

\begin{document}
\maketitle

% ========================================================
% ABSTRACT
% ========================================================

\begin{abstract}
This Companion develops the corrected interpretation of the coefficients
$C_1$–$C_4$ in RDCC. In v49.1, these coefficients are fixed theoretical factors
of order unity rather than fit parameters. Their values arise from the global
geometric structure of the bipartite Hilbert space and the ultraviolet-suppressed
dimension-8 operator $T^{(+)}_{\mu\nu}T^{(-)\mu\nu}$. This correction stabilizes
the Boltzmann hierarchy, the likelihood structure, and the parameter-estimation
pipeline, ensuring full EFT coherence and statistical consistency.
\end{abstract}

\tableofcontents
\newpage
% ========================================================
% 1 INTRODUCTION
% ========================================================

\section{Introduction}

The coefficients $C_1$–$C_4$ appear in several components of RDCC: the
Boltzmann hierarchy, the covariant Poisson equation, the likelihood structure,
and the survey interface. In earlier versions (up to v48.0), these coefficients
were treated as fit parameters. This interpretation is incorrect.

In v49.1, the coefficients $C_1$–$C_4$ are fixed theoretical factors of order
unity. Their values arise from the global geometric structure of the bipartite
Hilbert space and the ultraviolet-suppressed dimension-8 operator
\begin{equation}
\mathcal{O}_{\mathrm{int}} = T^{(+)}_{\mu\nu} T^{(-)\mu\nu}.
\end{equation}

This Companion develops the corrected interpretation of the $C$-coefficients
and explains their role in RDCC’s perturbation and likelihood structure.
% ========================================================
% 2 THEORETICAL ORIGIN OF THE C-COEFFICIENTS
% ========================================================

\section[Theoretical Origin]{Theoretical Origin}

The coefficients $C_1$–$C_4$ arise from the global geometric structure of the
bipartite Hilbert space. The global state is
\begin{equation}
\ket{\Psi} = \sum_n \lambda_n \ket{n}_{(+)} \otimes \ket{n}_{(-)},
\end{equation}
with reduced density matrices
\begin{align}
\rho_{(+)} &= \Tr_{(-)} \ket{\Psi}\bra{\Psi}, \\
\rho_{(-)} &= \Tr_{(+)} \ket{\Psi}\bra{\Psi}.
\end{align}

The inter-sector coupling is encoded in the dimension-8 operator
\begin{equation}
\mathcal{O}_{\mathrm{int}} = T^{(+)}_{\mu\nu} T^{(-)\mu\nu},
\end{equation}
which contributes to the covariant Poisson equation, the Boltzmann hierarchy,
and the likelihood structure.

The coefficients $C_1$–$C_4$ arise from geometric contractions of the global
metric and the energy-momentum tensors. They are therefore fixed theoretical
factors of order unity.
% ========================================================
% 3 FIXED VALUES AND ROLE IN RDCC
% ========================================================

\section[Fixed Values and Role]{Fixed Values and Role}

The coefficients satisfy
\begin{align}
C_1 &\sim \mathcal{O}(1), \\
C_2 &\sim \mathcal{O}(1), \\
C_3 &\sim \mathcal{O}(1), \\
C_4 &\sim \mathcal{O}(1).
\end{align}

They are not fit parameters. Their values are determined by the global
geometric structure and the ultraviolet suppression scale $M$.
\subsection[Role in the Boltzmann Hierarchy]{Role in the Boltzmann Hierarchy}

The coefficients appear in the covariant Poisson equation,
\begin{equation}
\nabla^2 \Phi - 3H(\dot{\Phi} + H\Phi)
=
4\pi G \left( \rho_{(+)} \delta_{(+)} + \rho_{(-)} \delta_{(-)} \right)
+
\frac{C_1}{M^4} \, \Xi,
\end{equation}
where $\Xi$ encodes the geometric contribution of the dimension-8 operator.
\subsection[Role in the Likelihood Structure]{Role in the Likelihood Structure}

The coefficients appear in the covariance matrix and determine the geometric
weighting of the residuals,
\begin{equation}
\mathcal{L}
=
\exp\left(
-\frac{1}{2}
\Delta\mathbf{d}^{\mathrm{T}}
\mathbf{C}^{-1}
\Delta\mathbf{d}
\right),
\end{equation}
where $\mathbf{C}$ contains the full cross-correlations between CMB, BAO, RSD
multipoles, supernovae, and gravitational-wave observables.
% ========================================================
% 4 CONCLUSION
% ========================================================

\section{Conclusion}

This Companion develops the corrected interpretation of the coefficients
$C_1$–$C_4$ in RDCC. In v49.1, these coefficients are fixed theoretical factors
of order unity rather than fit parameters. Their values arise from the global
geometric structure of the bipartite Hilbert space and the ultraviolet-suppressed
dimension-8 operator $T^{(+)}_{\mu\nu}T^{(-)\mu\nu}$.

These results stabilize the Boltzmann hierarchy (Companion VII), the likelihood
structure (Companion VIII), the parameter-estimation pipeline (Companion IX),
and the survey interface (Companion XII). Together, they ensure full EFT
coherence and statistical consistency.

\end{document}
